% SHARD-ID: RC-05-WEIL-LPS-CHANNELS
% SHARD-TITLE: Weil--LPS Channels: Exactly Ramanujan Quantum Expanders
% SHARD-SUMMARY: Builds the mixed-unitary channels obtained by feeding the LPS generators of PSL_2(F_p) through the Weil representation on l^2(F_p), and records the numerical checks of every construction step.
% SHARD-SUMMARY: The block channels are exactly Ramanujan by Harrow's transfer inequality plus Lubotzky-Phillips-Sarnak; their commuting joint spectra look like Hecke eigenvalues, which is left open.
% SHARD-KEYWORDS: Weil representation, LPS generators, quaternions, Ramanujan quantum expander, Hecke relation, joint spectrum, quantum Ihara zeta
% SHARD-NOTES: notes/weil-lps-channels.md
% SHARD-SCRIPTS: scripts/weil_lps.py, scripts/weil_lps_hashimoto.py

\section{Weil--LPS Channels: Exactly Ramanujan Quantum Expanders}
\label{sec:weil-lps}

This is the finite model of the Riemann channel of \cref{sec:riemann-channel}
with every ingredient present and every claim checkable: commuting dilations
indexed by primes, an arithmetic quotient supplying the expansion, and the
Weil representation as the Hilbert space. The construction was proposed in
the founding session's handoff and computed on 2026-09-11; everything below
was computed, and every construction step carries its own numerical check.
The object itself is \cref{def:weil-lps-channel}, its ingredients
\cref{def:lps-generators} and \cref{def:weil-representation}.

\subsection{The object}

Fix an odd prime $\pp$ and an LPS parameter $\qs\equiv1\pmod4$ with $\qs$ a
square mod $\pp$. The Weil intertwiner exists only for symplectic
(determinant one) matrices, and $\qs$ being a square is what lets the norm-$\qs$
quaternions land in $\SL(\Fp)$ after scaling. The admissible pairs computed
are these.

\begin{center}
\begin{tabular}{rrrl}
\toprule
$\pp$ & $\qs$ & $\lvert\PSL(\Fp)\rvert$ & Weil blocks $M_{(\pp+1)/2}$, $M_{(\pp-1)/2}$ \\
\midrule
$5$  & $29$     & $60$    & $M_3$, $M_2$ \\
$7$  & $29$     & $168$   & $M_4$, $M_3$ \\
$11$ & $5$      & $660$   & $M_6$, $M_5$ \\
$13$ & $17, 29$ & $1092$  & $M_7$, $M_6$ \\
$17$ & $13$     & $2448$  & $M_9$, $M_8$ \\
$19$ & $5, 17$  & $3420$  & $M_{10}$, $M_9$ \\
$29$ & $5, 13$  & $12180$ & $M_{15}$, $M_{14}$ \\
\bottomrule
\end{tabular}
\end{center}

\paragraph{Step 1, generators.} The $\qs+1$ integer quaternions
$\quat = a_0 + a_1 i + a_2 j + a_3 k$ of norm $\qs$ with $a_0>0$ odd and
$a_1,a_2,a_3$ even. Jacobi's four-square theorem gives exactly $\qs+1$ of
them; the enumeration was checked against the divisor sum
$\sum_{d\mid n}d$ for the norms $\qs$, $\qs^2$ and $\qs_1\qs_2$ used below. The same $\pp=5$ quaternions appear in
Appendix~C of \cite{iyer2026ramanujan}.

\paragraph{Step 2, images in $\SL(\Fp)$.} Split the Hamilton quaternions over
$\Fp$ by
\[
  \Isplit = \begin{pmatrix} x & y \\ y & -x\end{pmatrix},
  \qquad
  \Jsplit = \begin{pmatrix} 0 & 1 \\ -1 & 0\end{pmatrix},
  \qquad x^2+y^2+1\equiv0 \pmod \pp .
\]
Then $\quat\mapsto a_0 + a_1\Isplit + a_2\Jsplit + a_3\Isplit\Jsplit$ is a ring
homomorphism whose determinant is the quaternion norm. Dividing by a square
root $\sqrtq$ of $\qs$ mod $\pp$ gives $\galpha\in\SL(\Fp)$. Checked:
$\det\galpha = 1$, $g_{\bar{\quat}} = \galpha^{-1}$, and all $\qs+1$ images
distinct.

\paragraph{Step 3, the Weil representation, built and checked numerically.}
On $\wl^2(\Fp)$ take the clock and shift operators $\clock$, $\shift$ and the
Heisenberg--Weyl operators $\Hei = \shift^u\clock^v$. The Fourier transform
$\Four$, the chirp $\chirp$ and a dilation $\Dil$ each satisfy an
intertwining relation of the form
\[
  \Weil(g)\,\Hei\,\Weil(g)^\dagger \;\propto\; T(g(u,v)),
\]
with a symplectic $g$ that the code \emph{reads off} rather than assumes. A
breadth-first search over $\SL(\Fp)$ in these three generators assigns a
unitary $\Weil(g)$ to every group element, projectively; for $\pp=13$ the
table has the expected $\pp(\pp^2-1)=2184$ entries. The relation above was
then verified directly for each LPS generator $\galpha$, with error at most
$4\cdot10^{-14}$ in every case ($7.2\cdot10^{-15}$ for $(\pp,\qs)=(13,17)$,
$7.0\cdot10^{-15}$ for $(13,29)$). Since only $\Ad\Weil$ enters the channel,
the projective phases are irrelevant.

\paragraph{Step 4, the two irreducible blocks.} $\Weil(-1)$ is the parity
operator, so $\Weil$ splits into even and odd functions, of dimensions
$(\pp\pm1)/2$. Each block $\Wpm$ is irreducible and $-1$ acts on it as a
scalar, so it is a representation of $\PSL(\Fp)$ up to phase. This matters:
on the full $M_{\pp}$ the parity operator would be a second fixed point of the
channel, and Harrow's theorem needs irreducibility. The channels are
\[
  \PhiWL{\qs}^{\pm}(\dm)
  = \frac{1}{\qs+1}\sum_{\quat}\Wpm(\galpha)\,\dm\,\Wpm(\galpha)^\dagger
  \qquad\text{on } M_{(\pp\pm1)/2},
\]
that is, $\Dk = \qs+1$ unitary Kraus operators in the normalisation of
\cref{sec:conventions}. Block unitarity and Hilbert--Schmidt self-adjointness
of the superoperator were checked at $10^{-14}$ (for $(13,17)$: deviations
$2.7\cdot10^{-15}$ and $1.8\cdot10^{-15}$).

\subsection{Exact Ramanujan behaviour}

\begin{proposition}[Weil--LPS channels are exactly Ramanujan]
\label{prop:weil-lps-exactly-ramanujan}
\claimstatus{prop:weil-lps-exactly-ramanujan}{sketched}
For every admissible pair $(\pp;\qs)$, each block channel
$\PhiWL{\qs}^{\pm}$ is a Ramanujan quantum expander in the sense of
\cref{def:ramanujan-channel}: $\lamtwo(\PhiWL{\qs}^{\pm})\le\lamH$, equivalently
$(\qs+1)\,\lamtwo \le 2\sqrt{\qs}$ in adjacency units.
\end{proposition}

The argument is two lines. Harrow's transfer inequality
(\cref{cit:harrow-transfer}, \cite{harrow2008cayley}) says that for the channel
built from an irreducible representation of a group and a symmetric generating
set, the second singular value of the channel is at most the second singular
value of the transition matrix of the Cayley graph on that generating set. Here
the group is $\PSL(\Fp)$, the generating set is the $\qs+1$ images $\galpha$ of
\cref{def:lps-generators}, and the representation is $\Wpm$, irreducible by
Step 4; so $(\qs+1)\lamtwo(\PhiWL{\qs}^{\pm})$ is bounded by the largest modulus
of a nontrivial eigenvalue of the LPS Cayley graph $\Hecke{\qs}$. That graph is
Ramanujan, $\lvert\lambda\rvert\le2\sqrt{\qs}$, by Lubotzky, Phillips and Sarnak
\cite{lubotzky1988ramanujan}, and $2\sqrt{\qs} = 2\sqrt{\Dk-1}$ is exactly
Hastings' bound $\lamH$ in adjacency units \cite{hastings2007random}. The status
is \texttt{sketched} rather than \texttt{proved} because the Lubotzky--Phillips--Sarnak
theorem has no byte-verified local source in \texttt{refs/}, and because the
irreducibility of $\Wpm$ and the projective-phase bookkeeping of Step 4 are
verified numerically here rather than cited. Iyer, Jain, Jordan and Somma
\cite{iyer2026ramanujan} state the $\mathrm{SU}(2)$ version of the same
mechanism and call explicit Ramanujan quantum expanders a longstanding open
problem; see \cref{sec:prior-art}.

\begin{numerical}[Ramanujan, Harrow containment, Hecke, commutation]
\label{num:weil-lps-ramanujan}
\claimstatus{num:weil-lps-ramanujan}{numerical}
For $(\pp;\qs)$ in $\{(5;29),(7;29),(11;5),(13;17,29),(17;13),(19;5,17),(29;5,13)\}$
the channels $\PhiWL{\qs}^{\pm}$ on the even and odd Weil blocks satisfy the
Ramanujan bound of \cref{def:ramanujan-channel} and have a unique fixed point;
Harrow containment, the Hecke relation of \cref{def:hecke-relation} and
commutation across $\qs$ hold to $10^{-14}$.
\end{numerical}

Every entry of the table below is on the correct side of the bound
$2\sqrt{\qs}$, in adjacency units.

\begin{center}
\footnotesize
\begin{tabular}{rr rr rr rr r}
\toprule
& & \multicolumn{2}{c}{Cayley $\Hecke{\qs}$}
  & \multicolumn{2}{c}{$\PhiWL{\qs}^{+}$}
  & \multicolumn{2}{c}{$\PhiWL{\qs}^{-}$} & \\
\cmidrule(lr){3-4}\cmidrule(lr){5-6}\cmidrule(lr){7-8}
$\pp$ & $\qs$ & $\lamtwo$ & $\lambda_{\min}$
  & $(\qs{+}1)\lamtwo$ & $(\qs{+}1)\lambda_{\min}$
  & $(\qs{+}1)\lamtwo$ & $(\qs{+}1)\lambda_{\min}$ & $2\sqrt{\qs}$ \\
\midrule
$5$  & $29$ & $6.000$ & $-8.000$ & $6.000$ & $-2.000$ & $2.000$ & $2.000$ & $10.770$ \\
$7$  & $29$ & $8.000$ & $-6.000$ & $6.000$ & $-6.000$ & $2.000$ & $-6.000$ & $10.770$ \\
$11$ & $5$  & $4.372$ & $-4.024$ & $3.562$ & $-4.024$ & $3.406$ & $-4.024$ & $4.472$ \\
$13$ & $17$ & $6.928$ & $-7.851$ & $6.000$ & $-6.685$ & $5.555$ & $-5.555$ & $8.246$ \\
$13$ & $29$ & $9.932$ & $-9.208$ & $9.000$ & $-9.000$ & $9.000$ & $-9.000$ & $10.770$ \\
$17$ & $13$ & $6.106$ & $-7.090$ & $6.000$ & $-6.828$ & $6.000$ & $-6.828$ & $7.211$ \\
$19$ & $5$  & $4.275$ & $-4.274$ & $4.243$ & $-4.000$ & $4.243$ & $-4.000$ & $4.472$ \\
$19$ & $17$ & $7.671$ & $-7.638$ & $7.128$ & $-7.638$ & $7.128$ & $-7.638$ & $8.246$ \\
$29$ & $5$  & $4.442$ & $-4.411$ & $4.442$ & $-4.049$ & $4.442$ & $-4.049$ & $4.472$ \\
$29$ & $13$ & $6.765$ & $-6.949$ & $6.765$ & $-6.912$ & $6.579$ & $-6.912$ & $7.211$ \\
\bottomrule
\end{tabular}
\end{center}

The further checks, all passed in every case, are these.

\begin{itemize}
  \item \textbf{Harrow containment.} Every channel eigenvalue, times $\qs+1$,
  lies in the spectrum of the $\PSL(\Fp)$ Cayley graph to $5\cdot10^{-7}$
  (for $(13,17)$: maximum distance $4.4\cdot10^{-7}$ on the even block,
  $3.1\cdot10^{-7}$ on the odd). The graph spectrum was computed in double
  precision on up to $3420$ vertices; for $\pp=29$ only the extreme
  eigenvalues were computed, by Lanczos, which is why this is a containment
  check and not an equality check.
  \item \textbf{Unique fixed point.} The eigenvalue $1$ of each block channel
  has multiplicity one, so the channels are quantum expanders in the sense of
  \cref{def:quantum-expander} and not merely bounded.
  \item \textbf{Hecke relation.} $\Hecke{\qs}^2 - \qs\,1 = \Hecke{\qs^2}$ on
  $\PSL(\Fp)$, with $\Hecke{\qs^2}$ built from the $\qs^2+\qs+1$ normalised
  quaternions of norm $\qs^2$ ($307$ for $\qs=17$, $871$ for $\qs=29$).
  Residual exactly $0$. This is the tree relation ``two steps without
  backtracking''.
  \item \textbf{Commutation across $\qs$.} For
  $(\pp;\qs_1,\qs_2) = (13;17,29)$, $(19;5,17)$, $(29;5,13)$ the commutator
  $\lVert[\PhiWL{\qs_1}^{\pm},\PhiWL{\qs_2}^{\pm}]\rVert$ is at most
  $3\cdot10^{-15}$ (for $\pp=13$: $5.4\cdot10^{-16}$ even,
  $4.5\cdot10^{-16}$ odd), and $\Hecke{\qs_1}\Hecke{\qs_2} =
  \Hecke{\qs_1\qs_2}$ exactly, with the $(\qs_1+1)(\qs_2+1) = 540$ quaternions
  of norm $\qs_1\qs_2$.
\end{itemize}

\subsection{The quantum Ihara zeta of these channels}

\begin{numerical}[Edge spectrum and Ihara--Bass for $(\pp,\qs)=(13,17)$]
\label{num:weil-lps-rh}
\claimstatus{num:weil-lps-rh}{numerical}
For $(13,17)$, odd block, the non-backtracking edge superoperator $\Hash$ on
$M_6\otimes\mathbb C^{18}$, of dimension $648$, has $\lvert\edgeev\rvert$ in
$\{17,\ \sqrt{17},\ 1\}$ with multiplicities $1$, $70$, $577$, and satisfies
the Ihara--Bass identity to $10^{-14}$.
\end{numerical}

The edge superoperator of \cref{def:quantum-hashimoto} was built directly, and
the Ihara--Bass factorisation of \cref{def:ihara-bass},
\[
  \det(1-\uz\Hash) \;=\; (1-\uz^2)^{\,n^2(\Dk-2)/2}\,
  \det\!\big(1 - \uz\,\Dk\,\PhiWL{\qs}^{-} + \qs\,\uz^2\big),
  \qquad n=6,\ \Dk=18,
\]
holds with relative residuals $2.8\cdot10^{-15}$ at $\uz=0.11$,
$2.3\cdot10^{-15}$ at $\uz=0.2+0.07i$ and $1.5\cdot10^{-14}$ at $\uz=-0.23$.
The multiplicities $1+70+577 = 648$ exhaust the edge space, and
$\sqrt{\qs} = 4.12311$. So the quantum Ihara zeta $\zetaq$ of this channel
satisfies the RH analogue of \cref{def:rh-analogue} on the nose: every
nontrivial pole sits on the critical circle (the ungraded quantum Ihara zeta has poles only;
wording corrected 2026-09-15 after the prover lane of \cref{sec:graded-ramanujan}). For every other case the same
conclusion follows from the channel spectrum through the quadratic relation
$\edgeev^2 - \lambda\,\edgeev + \qs = 0$ that links an edge eigenvalue to an
adjacency eigenvalue $\lambda$, without building $\Hash$.

\subsection{What the joint spectrum looks like}

For $(13;17,29)$ the even block has $19$ distinct joint eigenvalue pairs
$(a_{17}, a_{29})$ of the commuting pair $\PhiWL{17}^{+}$, $\PhiWL{29}^{+}$.
Examples, in adjacency units.

\begin{center}
\begin{tabular}{lll}
\toprule
$(a_{17}, a_{29})$ & multiplicity & remark \\
\midrule
$(18, 30)$ & $1$ & trivial: $\qs+1$ \\
$(6, 2)$ & $2$ & integers \\
$(3,-1)$, $(3,3)$, $(3,9)$ & $3$, $3$, $2$ & integers \\
$(-5,-5)$ & $6$ & integers \\
$\big(\tfrac{3+\sqrt{17}}{2},\ -1-\sqrt{17}\big)$ & $3$ & quadratic, field $\mathbb Q(\sqrt{17})$ \\
$(5.6847, -2)$, $(5.5549, -0.3344)$, $(1.7730, 5.1119)$, \dots & $3$ each & higher degree \\
\bottomrule
\end{tabular}
\end{center}

That the eigenvalues are algebraic integers is not a test: the Cayley
adjacency matrix $\Hecke{\qs}$ is an integer matrix by construction. What is
visible is the structure a Hecke identification would predict, namely joint
eigenvalue pairs, small-degree algebraic integers, multiplicities equal to the
multiplicity of the corresponding irreducible representation of $\PSL(\Fp)$ in
$\Wpm\otimes\overline{\Wpm}$, and every pair inside the box
$[-2\sqrt{\qs_1},2\sqrt{\qs_1}]\times[-2\sqrt{\qs_2},2\sqrt{\qs_2}]$.

\begin{conjecture}[Hecke identification of the Weil--LPS spectra]
\label{conj:weil-lps-hecke}
\claimstatus{conj:weil-lps-hecke}{open}
The joint spectra of the commuting Weil--LPS channels $\PhiWL{\qs}^{\pm}$ are
Hecke eigenvalues $a_{\qs}(f)$ of weight-$2$ forms attached to the quaternion
algebra ramified at $\{2,\infty\}$, with level determined by $\pp$, via
Jacquet--Langlands.
\end{conjecture}

The identification was \emph{not} done here; it needs LMFDB and the
Jacquet--Langlands bookkeeping.

\subsection{Reading}

\begin{center}
\begin{tabular}{ll}
\toprule
Riemann channel & Weil--LPS channel at $\pp$ \\
\midrule
commuting dilations by primes & commuting Hecke channels $\PhiWL{\qs}$, one per LPS prime $\qs$ \\
arithmetic quotient supplying expansion & $\PSL(\Fp)$ as quotient of the $(\qs{+}1)$-regular tree \\
Weil representation & $\Wpm$ on $\wl^2(\Fp)$, even and odd \\
joint spectrum $=$ zeros & joint spectrum $=$ Hecke eigenvalues \\
RH: all modes at $\re = 1/2$ & Ramanujan: all modes at $\lvert\edgeev\rvert = \sqrt{\qs}$ \\
source of the bound: unknown & source of the bound: Deligne, via LPS \\
\bottomrule
\end{tabular}
\end{center}

The last row is the honest one. These channels are exactly Ramanujan because
Deligne proved the Weil conjectures (\cref{sec:deligne}), not because they are
channels. What the computation establishes is that the object proposed in the
founding session's handoff exists, is cheap to build, is self-checking at every
step, and lands on the parent campaign's Hilbert space with the parent
campaign's representation.

\subsection{Not established}

\begin{itemize}
  \item The identification of the joint spectrum with specific modular forms
  in LMFDB (\cref{conj:weil-lps-hecke}).
  \item Any statement about $\pp\to\infty$, or about assembling the
  $\PhiWL{\qs}$ over all $\pp$ into an object whose rings are the primes; that
  is the parent repository's global question in a new guise, recorded as
  \cref{conj:assemble-over-p} in \cref{sec:open}.
  \item Anything about the Riemann hypothesis itself.
\end{itemize}
